\section{Conclusión}

Los resultados experimentales sobre los microdatos de la ENAHO 2024--2025 permiten responder de forma concluyente la pregunta de investigación: en el contexto urbano e informal de Lima Metropolitana y el Callao, \textbf{el factor determinante en la detección de la pobreza monetaria es la riqueza multidimensional de la representación del hogar ($\mathbb{R}^{28}$) y no la sofisticación algorítmica del clasificador}. Mientras que pasar de indicadores físicos de vivienda (Módulo 01) a la representación conjunta con demografía, educación y empleo formal/informal incrementa la PR-AUC fuera de tiempo entre $+0.210$ y $+0.232$ y reduce el error de exclusión a cobertura $B=20\%$ en cerca de $18$ puntos porcentuales (de $0.582\text{--}0.591$ a $0.393\text{--}0.409$, corroborando $H_1$ al 95\% de confianza), la diferencia de error de exclusión entre una línea base econométrica bien regularizada (\texttt{PMT\_MCO} con Ridge), una regresión logística ElasticNet, tres librerías de \textit{gradient boosting} (\texttt{XGBoost}, \texttt{CatBoost}, \texttt{LightGBM}) y el meta-ensamble ganador (\texttt{Stacking\_Ensemble}) no supera $1.6$ puntos porcentuales ($H_2$ corroborada al 95\% para \texttt{Stacking}, \texttt{XGBoost} y \texttt{CatBoost}). Asimismo, al corregir la imputación espacial mediante la cascada libre de fuga \texttt{conglomerado} $\rightarrow$ \texttt{ubigeo} $\rightarrow$ mediana global, el umbral $\tau_{2024}$ congelado en 2024 logra simultáneamente en 2025 una exhaustividad $\ge 0.70$ y una precisión $\ge 0.45$ en los siete modelos principales ($H_3$), con una caída interanual de PR-AUC inferior a $0.026$ puntos en estimación puntual ($H_4$), aunque con intervalos de confianza al 95\% más amplios por la variabilidad muestral entre conglomerados. En suma, combinar un meta-ensamble regularizado con atribuciones aditivas TreeSHAP ofrece al sistema de focalización social una herramienta precisa, estable entre años y transparente para auditar cada decisión a nivel de hogar.
